% Copyright (c) 2026 Eugene Joseph M. Ragasa.
% SPDX-License-Identifier: Apache-2.0
% Author: Eugene Joseph M. Ragasa
% Formal proof completed: 2026-10-20
\documentclass[11pt]{article}
\usepackage[T1]{fontenc}\usepackage{amsmath,amssymb}\usepackage[margin=1in]{geometry}\usepackage[hidelinks]{hyperref}
\title{Proof Dossier: Hermitian Rank-Two Pauli Representation}\author{Eugene Joseph M. Ragasa}\date{Formal proof completed: 2026-10-20}
\begin{document}\maketitle
\section{Identity and status}
ID \texttt{BG-FOUNDATION-THM-013}; status \texttt{proved}. Lean declaration:
\par\noindent{\tiny\path{ProjectKoios.BandGauge.Foundations.}}\\
{\tiny\path{eq_rankTwoHermitianMatrix_of_conjTranspose_eq_self}}.
\section{Source and contribution boundary}
Source-neutral supporting algebra; relation \texttt{background-only} to standard finite
Hermitian matrix theory through pinned Mathlib. No publication-source or novelty claim.
\section{Statement and proof}
Every Hermitian complex $2$-by-$2$ matrix is
\[
\begin{pmatrix}t+z&x-iy\\x+iy&t-z\end{pmatrix},
\]
where $t$ is half the sum of the real diagonal entries, $z$ is half their difference,
$x=\operatorname{Re}M_{01}$, and $y=-\operatorname{Im}M_{01}$. Hermiticity makes both
diagonal entries real and identifies the lower-left entry with the conjugate upper-right
entry; entrywise substitution completes the proof.
\section{Boundary}
This representation connects arbitrary Hermitian rank-two matrices to the coordinate
perturbation theorem. It does not certify numerical Hermiticity or eigensolver output.
Kernel axioms are \texttt{propext}, \texttt{Classical.choice}, and \texttt{Quot.sound}; no admissions. Checked with Lean \texttt{4.35.0} and Mathlib revision \texttt{56a005855ea17e0b0958f73ec801ae140b0a1679}. AI-assisted and kernel-checked. No independent expert review is recorded.
\end{document}
